\section*{Цель работы}
Изучение характеристик и принципов функционирования аналоговых электронных ключей, реализованных на различных полупроводниковых приборах (диодах, биполярных и полевых транзисторах, оптронах).

\section*{Задачи}
\begin{enumerate}
    \item Снятие статических характеристик (напряжения в открытом и закрытом состояниях).
    \item Определение коэффициентов коммутации и передачи.
    \item Исследование амплитудно-частотных характеристик (АЧХ) ключей.
    \item Оценка быстродействия и определение пороговых напряжений.
    \item Измерение входного сопротивления по управляющему входу.
\end{enumerate}

\section*{Краткие теоретические сведения}
Аналоговый электронный ключ предназначен для коммутации сигналов от источника к нагрузке. Основные параметры:
\begin{itemize}
    \item \textit{Коэффициент коммутации} $K_k$: отношение напряжений на выходе в открытом и закрытом состояниях. Для последовательного ключа:
    \begin{equation*}
        K_k = \frac{U_{вых.откр}}{U_{вых.закр}}
    \end{equation*}
    \item \textit{Коэффициент передачи} $K_п$:
    \begin{equation*}
        K_п = \frac{U_{вых.откр}}{U_{вх}}
    \end{equation*}
    \item \textit{Быстродействие}: время перехода между уровнями:
    \begin{equation*}
        0.1 U_{вых.max}~и~0.9 U_{вых.max}
    \end{equation*}
\end{itemize}

\section*{Инструкция по выполнению работы}

\textit{Перед началом работы убедитесь, что все приборы заземлены, а напряжения питания на макете соответствуют требуемым.}

\subsection*{1. Исследование диодного ключа}
\begin{enumerate}
    \item \textit{Подготовка схемы} $\rightarrow$ Соберите схему согласно \cref{fig:circuit_diode}. Установите напряжения питания $\pm 12$~В. $\rightarrow$ Важно соблюдать полярность подключения источников.
    \item \textit{Настройка сигналов} $\rightarrow$ Подайте на вход гармонический сигнал $U_{вх} = 0,5$~В ($f = 1$~кГц). На управляющий вход подайте меандр $U_{упр} = 10$~В, длительностью $5$~мс. $\rightarrow$ Контролируйте форму сигналов по осциллографу.
    \item \textit{Статические измерения} $\rightarrow$ Измерьте $U_{вых}$ для двух состояний ключа. Занесите данные в \cref{tab:static_params}. $\rightarrow$ Используйте режим курсорных измерений осциллографа для точности.
    \item \textit{Снятие АЧХ} $\rightarrow$ При неизменном $U_{вх}$ варьируйте частоту $f$ от $1$~кГц до $1$~МГц. Фиксируйте $U_{вых}$ в \cref{tab:frequency_response}. $\rightarrow$ Следите, чтобы амплитуда генератора не «плавала» при смене диапазона.
    \item \textit{Оценка быстродействия} $\rightarrow$ Установите $f = 1$~кГц. Измерьте время переключения при $U_{упр} = 10$~В и $U_{упр} = 2$~В. Данные занесите в \cref{tab:speed_params}.
    \item \textit{Пороговое напряжение} $\rightarrow$ Плавно изменяя амплитуду $U_{упр}$, определите момент резкого изменения сопротивления ключа.
    \item \textit{Входное сопротивление} $\rightarrow$ Подключите магазин сопротивлений между генератором импульсов и управляющим входом. Увеличивайте сопротивление до тех пор, пока амплитуда импульса на входе макета не уменьшится в 2 раза. Запишите значение сопротивления магазина.
\end{enumerate}

\subsection*{2. Исследование транзисторных и оптронных ключей}
\begin{enumerate}
    \item \textit{Последовательность действий} $\rightarrow$ Повторите шаги п. 1 (б--и) для схем на биполярном транзисторе (\cref{fig:circuit_bipolar}), полевом транзисторе (\cref{fig:circuit_fet}) и оптроне (\cref{fig:circuit_opto}).
    \item \textit{Особенности для оптрона} $\rightarrow$ Для схемы на \cref{fig:circuit_opto} установите входное напряжение $U_{вх} = 1$~В. $\rightarrow$ Обратите внимание на гальваническую развязку цепей.
\end{enumerate}

\begin{figure}[H]
    \centering
    \includegraphics[width=0.4\textwidth]{figs/circuit_diode.pdf}
    \caption{Схема последовательного диодного ключа}
    \label{fig:circuit_diode}
\end{figure}

\begin{figure}[H]
    \centering
    \includegraphics[width=0.4\textwidth]{figs/circuit_bipolar.pdf}
    \caption{Схема параллельного ключа на биполярном транзисторе}
    \label{fig:circuit_bipolar}
\end{figure}

\begin{figure}[H]
    \centering
    \includegraphics[width=0.4\textwidth]{figs/circuit_fet.pdf}
    \caption{Схема ключа на полевых транзисторах}
    \label{fig:circuit_fet}
\end{figure}

\begin{figure}[H]
    \centering
    \includegraphics[width=0.4\textwidth]{figs/circuit_opto.pdf}
    \caption{Схема ключа на оптронах}
    \label{fig:circuit_opto}
\end{figure}

% Соответствие рисунков методическим материалам:
% fig:circuit_diode   <-> Рис. 9.3
% fig:circuit_bipolar <-> Рис. 9.4
% fig:circuit_fet     <-> Рис. 9.5
% fig:circuit_opto    <-> Рис. 9.6